\section{Contributions}
This thesis aims to transform the given type one algorithm and optimize it for streaming data. One of the requirements is that the Principal Component Analysis should be extended to be calculated on an sliding window which is adjustable in size. This means that new data points can be added to the time series while reusing the already calculated principal component analysis. Instead of recomputing the covariance matrix, eigenvectors, and eigenvalues from scratch, the algorithm should efficiently update these values using the previously calculated values. Another main contribution is the integration of a sliding window mechanism which bounds memory usage and ensures the responsiveness of the model to recent data patterns. Only the last $k$-elements from the time series should be part of the equation additionally partial results should be reused. Older points are pushed out of the sliding window ensuring that the algorithm can adapt to changing signal patterns while holding a constant memory usage. \\
The methodology section explains the reasoning behind the choosing of the individual parts of the tech stack. This includes programming language and frameworks. Finally, the optimized implementation is evaluated by analyzing how the parameters window size, delay, stride and embedding dimensions impact the visual representation of the resulting fingerprint. For easy integration the result will be a small library with an application programming interface (API) which lets the user easily integrate principal component analysis in  software projects.
